\documentclass{article}
\usepackage{amsmath,amssymb,amsthm}
\newcommand{\Om}{\Omega}
\newtheorem{lemma}{Lemma}
\begin{document}
\section{Decay for the heat equation}
Let $\Om\subset\mathbb R^N$ be a bounded smooth domain and let $\mu>0$
satisfy
\begin{equation}\label{eq:poincare-input}
 \int_{\Om}|\nabla v|^2\,dx\ge\mu\int_{\Om}|v|^2\,dx
 \qquad\text{for every }v\in H_0^1(\Om).
\end{equation}
This inequality is an assumption; no optimality of $\mu$ is needed.
\begin{lemma}\label{lem:heat-decay}
Let $T>0$ and let $u$ be a smooth real-valued solution of
$u_t=\Delta u$ on $[0,T]\times\Om$, with $u=0$ on
$[0,T]\times\partial\Om$. Then
\begin{equation}\label{eq:heat-conclusion}
 \|u(t)\|_{L^2(\Om)}\le e^{-\mu t}\|u(0)\|_{L^2(\Om)}
 \qquad(0\le t\le T).
\end{equation}
\end{lemma}
\begin{proof}
Define the squared $L^2$ energy by
\[
 E(t)=\int_{\Om}u(t,x)^2\,dx=\|u(t)\|_{L^2(\Om)}^2,
 \qquad 0\le t\le T.
\]
Smoothness allows us to differentiate under the integral sign. Using the
heat equation and integrating by parts, we obtain
\[
 E'(t)=2\int_{\Om}u u_t\,dx
      =2\int_{\Om}u\Delta u\,dx
      =-2\int_{\Om}|\nabla u|^2\,dx.
\]
The boundary term vanishes because $u=0$ on $\partial\Om$.
For each time $t$, the function $u(t,\cdot)$ belongs to $H_0^1(\Om)$, so
\eqref{eq:poincare-input} yields
\[
 E'(t)\le -2\mu E(t).
\]
Multiplying by the integrating factor $e^{2\mu t}$ gives
\[
 \frac{d}{dt}\bigl(e^{2\mu t}E(t)\bigr)
 =e^{2\mu t}\bigl(E'(t)+2\mu E(t)\bigr)\le0.
\]
Integrating from $0$ to $t$ therefore gives
$E(t)\le e^{-2\mu t}E(0)$. Taking square roots proves
\eqref{eq:heat-conclusion}. This last step also covers the case $E(0)=0$;
no division by the energy is needed.
\end{proof}
\end{document}
